\subsection{Token-Count-Matched Random Pruning}
\label{app:token_matched}

The message-count-matched random baseline in Table~\ref{tab:main} prunes the same
\emph{number} of messages as ack-only but not the same volume of content.
Because acknowledgments are short (Qwen: $\sim$12 est.\ tokens vs.\ $\sim$64 overall mean;
Llama: $\sim$80 vs.\ $\sim$110), count-matched random pruning removes substantially more
tokens. We test whether this token asymmetry explains the accuracy gap by constraining
random pruning to match ack-only's \emph{token} budget rather than its message count.

\paragraph{Methodology.} For each backbone, we compute the total estimated tokens
in ack-labeled messages. We then perform greedy-knapsack random sampling (10~seeds):
messages are shuffled and greedily added to the prune set until the ack token budget
is exhausted. Accuracy risk is measured via offline neutralization (fraction of correct
tasks whose final answer changes when pruned messages are removed from context).

\begin{table}[h]
\centering
\caption{Offline accuracy risk (\%) under three pruning regimes.
\emph{Count-matched} matches ack message count (removes $k{\times}$ more tokens);
\emph{token-matched} constrains to the same token budget.
Risk = fraction of correct tasks whose final answer changes upon message removal
(10~seeds, mean $\pm$ SD for random).}
\label{app:tab:token_matched}
\small
\begin{tabular}{lcccc}
\toprule
\textbf{Backbone} & \textbf{Ack-only} & \textbf{Count-matched} & \textbf{Token-matched} & \textbf{Token ratio} \\
 & risk (\%) & risk (\%) & risk (\%) & (count/ack) \\
\midrule
Qwen & 2.6 & 1.2$\pm$0.3 & 0.2$\pm$0.2 & 5.51$\times$ \\
Llama & 1.8 & 1.2$\pm$0.6 & 0.8$\pm$0.6 & 1.38$\times$ \\
Gemma & 0.0 & 0.0$\pm$0.0 & 0.0$\pm$0.0 & 8.14$\times$ \\
\bottomrule
\end{tabular}
\end{table}

\paragraph{Results.}
For Qwen, count-matched random removes 5.51$\times$
more tokens than ack-only (9\% of total context).
When constrained to the same token budget (8235 est.\ tokens),
random pruning touches only $\sim$138$ messages
(vs.\ 719 for ack-only) and achieves
0.2\% risk---lower than ack-only's
2.6\%.
This occurs because ack-only targets a specific message class that occasionally
contains answer-relevant content (2.6\% of ack-suppressed tasks change their final answer),
whereas distributing the same small token budget randomly across many partial messages
is unlikely to hit any single answer-critical message.
For Llama, the gap narrows from 1.38$\times$ token
overshoot to matched budgets, with token-matched random risk at
0.8\%
vs.\ ack-only's 1.8\%.

\paragraph{Interpretation.}
The offline token-matched analysis suggests that ack-only's safety advantage over
count-matched random in the e2e experiments (Table~\ref{tab:main}: Qwen ack 92\%
vs.\ random 87.3\%) is primarily driven by the $\sim$5$\times$ token asymmetry rather
than by the semantic targeting of acknowledgments.
However, this finding must be interpreted cautiously: (1)~offline risk simulation
assumes message independence, while e2e pruning involves causal cascade effects;
(2)~ack-only's value lies not in raw accuracy advantage over token-matched random,
but in its \emph{interpretability}---it targets a semantically coherent message class
with low information content, whereas random pruning at matched token budgets would
arbitrarily truncate or skip substantive content in unpredictable ways.